\chapter{Stevens-Johnson Syndrome and Toxic Epidermal Necrolysis}
\index{Stevens-Johnson Syndrome and Toxic Epidermal Necrolysis}
\label{sec:Stevens-Johnson Syndrome and Toxic Epidermal Necrolysis}

\section{Overview}
Stevens-Johnson syndrome and toxic epidermal necrolysis are severe, life-threatening mucocutaneous reactions characterized by extensive keratinocyte apoptosis and epidermal detachment, with SJS affecting $<$10\% of body surface area, SJS/TEN overlap affecting 10--30\%, and TEN affecting $>$30\% BSA. This genetic disorder results from specific gene mutations or chromosomal abnormalities. It may be present at birth or manifest later in life depending on the underlying mechanism. Genetic disorders result from abnormalities in an individual's DNA, ranging from single-gene mutations to chromosomal abnormalities. These conditions may be inherited from parents or arise from new mutations. Pattern of inheritance depends on whether the mutation is dominant, recessive, or X-linked. Genetic counseling helps families understand risks and make informed decisions.
\section{Causes}
The most common cause is drug hypersensitivity (anticonvulsants [carbamazepine, lamotrigine, phenytoin], allopurinol, sulfonamides, NSAIDs, antibiotics). Genetic disorders result from mutations in specific genes that encode proteins essential for normal cellular function. The pattern of inheritance depends on whether the mutation is dominant, recessive, or X-linked.   
\section{Risk Factors}

\begin{itemize}[noitemsep]
  \item Genetic predisposition and family history
  \item Advancing age
  \item Lifestyle factors (diet, exercise, smoking, alcohol)
  \item Environmental exposures
  \item Underlying medical conditions
  \item Medication use
\end{itemize}
\section{Signs and Symptoms}

\subsection{Common symptoms}
\begin{itemize}[noitemsep]
  \item \item You may experience prodromal fever and malaise, followed by painful, dusky red to purpuric macules that coalesce and progress to epidermal detachment with flaccid blisters and erosion, with mucosal involvement (oral, ocular, genital) present in over 90\%.

  \item Pain or discomfort in the affected area
  \end{itemize}

\subsection{Emergency warning signs}
\begin{itemize}[noitemsep]
  \item Severe or worsening symptoms of stevens-johnson syndrome and toxic epidermal necrolysis require immediate medical evaluation
\end{itemize}
\section{Progression and Stages}
The condition may progress through different stages of severity over time. Early stages often involve mild or subtle symptoms that may be overlooked. Moderate stages involve more noticeable symptoms affecting daily function. Advanced stages may involve significant impairment and complications.
\section{Complications}
\begin{itemize}[noitemsep]
  \item Disease progression leading to worsening symptoms
  \item Secondary infections or injuries
  \item Side effects of treatment
  \item Reduced quality of life
  \item Emotional and psychological impact
\end{itemize}
\section{Diagnosis}
Doctors usually diagnose this based on your symptoms and a physical exam. Management requires burn unit or ICU care: fluid resuscitation, wound care, nutritional support, ophthalmological consultation, pain management, and withdrawal of the offending drug, with cyclosporine and IVIG having shown some benefit.

\begin{itemize}[noitemsep]
  \item Comprehensive medical history and physical examination
  \item Laboratory tests to assess organ function
  \item Imaging studies to visualize affected structures
  \item Specialized testing depending on the affected system
  \item Consultation with relevant specialists
\end{itemize}
\section{Prevention}
\begin{itemize}[noitemsep]
  \item Maintain a healthy lifestyle with balanced diet and exercise
  \item Avoid known risk factors
  \item Attend regular medical check-ups for early detection
  \item Follow recommended screening guidelines
\end{itemize}
\section{Treatment Options}
Symptomatic management and supportive care Enzyme replacement therapy for certain conditions Gene therapy (emerging for select conditions) Dietary modifications (e.g., phenylketonuria) Regular monitoring for associated complications Physical and occupational therapy Genetic counseling for family planning Participation in clinical trials for new therapies

\begin{itemize}[noitemsep]
  \item Medications to manage symptoms and address underlying mechanisms
  \item Lifestyle modifications and self-care measures
  \item Physical therapy and rehabilitation
  \item Surgical intervention when indicated
  \item Regular monitoring and follow-up care
\end{itemize}
\section{Living with It}
\begin{itemize}[noitemsep]
  \item Follow your treatment plan as prescribed
  \item Maintain regular follow-up with your healthcare provider
  \item Make necessary lifestyle adjustments
  \item Seek support from family, friends, or support groups
  \item Stay informed about your condition
\end{itemize}
\section{Outlook}
outlook is guarded, with mortality of 10--30\% for SJS and up to 50\% for TEN.
